\begin{table}[t]
\centering
\caption{Effect of the 2020 Equal Pay Act on the Regular/Non-Regular Wage Gap}
\label{tab:main}
\small
\begin{tabular}{lccccc}
\toprule
 & Gap ratio & Log gap & Non-reg. wage & Regular wage & Industry \\
 & (1) & (2) & (3) & (4) & (5) \\
\midrule
Post-reform & 1.20*** & 0.020*** & -3.54*** & -12.91*** & \\
 & (0.07) & (0.001) & (0.06) & (0.65) & \\
Intensity $\times$ Post & & & & & 0.08 \\
 & & & & & (0.25) \\
\midrule
Observations & 33 & 33 & 33 & 33 & 132 \\
Within $R^2$ & 0.036 & 0.040 & 0.033 & 0.141 & 0.000 \\
Panel \& year FE & Yes & Yes & Yes & Yes & Yes \\
\bottomrule
\multicolumn{6}{p{0.98\textwidth}}{\footnotesize \textit{Notes:} Columns 1--4 use the firm-size panel (3 sizes $\times$ 11 years = 33 obs). Column 5 uses the balanced industry panel (12 industries $\times$ 11 years = 132 obs). The dependent variable in Column 1 is the wage gap ratio (non-regular/regular $\times$ 100); higher values indicate a narrower gap. Column 2 uses the log wage gap. Columns 3--4 decompose the effect into non-regular and regular wage level changes (thousands of yen). Column 5 tests whether industries with larger pre-reform gaps (standardized treatment intensity) experienced larger post-reform changes. All specifications include panel and year fixed effects. Standard errors clustered by firm size (Columns 1--4) or industry (Column 5) in parentheses. * $p<0.1$, ** $p<0.05$, *** $p<0.01$.} \\
\end{tabular}
\end{table}
